% Source: https://artofproblemsolving.com/community/c3h1392645_assorted_algebra_problems
% Stable identity: aops-topic-1392645-q03
From the first equation $y=2x^2-4$. Substituting into the second:
\[
2(2x^2-4)^2=x+7\iff 8x^4-32x^2-x+25=0 .
\]
Each real root $x$ of this quartic gives exactly one real solution $(x,2x^2-4)$ of the system, and conversely.

\textbf{(a)} Let $f(x)=8x^4-32x^2-x+25$. Then
\[
f(-2)=27,\quad f\left(-\dfrac{3}{2}\right)=-5,\quad f(-1)=2,\quad f(1)=0,\quad f\left(\dfrac{3}{2}\right)=-8,\quad f(2)=23 .
\]
By the intermediate value theorem $f$ has a root in each of $\left(-2,-\dfrac{3}{2}\right)$, $\left(-\dfrac{3}{2},-1\right)$, $\left(\dfrac{3}{2},2\right)$, and $x=1$ is a fourth root. A quartic has at most four roots, so there are exactly four real solutions.

\textbf{(b)} Dividing by $x-1$ gives
\[
f(x)=(x-1)(8x^3+8x^2-24x-25).
\]
The other three roots are the roots of the cubic $g(x)=8x^3+8x^2-24x-25$, one in each of the three intervals above. By the rational root theorem a rational root of $g$ has the form $\pm\dfrac{p}{q}$ with $p\mid25$ and $q\mid8$. The candidates are $\pm1,\pm5,\pm25$ divided by $1,2,4$ or $8$, and the only one lying in the three intervals is $-\dfrac{5}{4}$. It is not a root:
\[
g\left(-\dfrac{5}{4}\right)=-\dfrac{125}{8}+\dfrac{25}{2}+30-25=\dfrac{15}{8}\ne0 .
\]
So the three other real solutions are irrational, and the only rational solution is $x=1$, $y=2\cdot1-4=-2$:
\[
\boxed{(x,y)=(1,-2)}.
\]
